\documentclass[UTF8,11pt]{ctexart}
\usepackage[a4paper,margin=22mm]{geometry}
\usepackage{amsmath,amssymb,bm,graphicx,adjustbox,fvextra,longtable,array}
\xeCJKDeclareCharClass{CJK}{"2160 -> "217F, "2460 -> "249B, "2200 -> "22FF}
\newcommand{\fitmath}[1]{\begin{adjustbox}{max width=\linewidth}$\displaystyle #1$\end{adjustbox}}
\setlength{\parindent}{0pt}\setlength{\parskip}{.65em}\renewcommand{\arraystretch}{1.3}
\sloppy\allowdisplaybreaks
\begin{document}
\section*{2022 年考研政治真题}
\subsection*{2022年全国研究生入学考试政治试题}

\subsection*{一、单项选择题：1～16小题，每小题1分，共16分。下列每题给出的四个选项中，只有一个选项是符合题目要求的。请在答题卡上将所选项的字母涂黑。}

1. 中国共产党坚持马克思主义基本原理，坚持实事求是，从中国实际出发， 洞察时代大势， 把握历史动力，进行艰辛探索,不断推进马克思主义中国化时代化，指导中国人民不断推进伟大社会革命。习近平总书记指出: “中国共产党为什么能，中国特色社会主义为什么好，归根到底是因为马克思主义行! ”马克思主义之所以“行”，根本原因在于


\begin{center}\begin{tabular}{>{\raggedright\arraybackslash}p{\dimexpr(\linewidth-4\tabcolsep)/2\relax}>{\raggedright\arraybackslash}p{\dimexpr(\linewidth-4\tabcolsep)/2\relax}}
\textbf{A.} 马克思主义具有鲜明的政治立场 & \textbf{B.} 马克思主义具有自觉的历史担当\\[.35em]
\textbf{C.} 马克思主义是科学的世界观和方法论 & \textbf{D.} 马克思主义是无产阶级政党自我革命的武器\\[.35em]\end{tabular}\end{center}

答案：C


2. 党的十八大以来，我国从中西部 22 个省份有劳动能力的建档立卡贫困人口中选聘了 110.2 万名生产护林员, 走出了一条生态补偿脱贫的新路子， 实现了生态保护和脱贫增收“双赢”，充分体现了“人不负青山，青山定不负人”的深刻哲理，“人不负青山，青山定不负人”表明


\begin{center}\begin{tabular}{>{\raggedright\arraybackslash}p{\dimexpr(\linewidth-4\tabcolsep)/2\relax}>{\raggedright\arraybackslash}p{\dimexpr(\linewidth-4\tabcolsep)/2\relax}}
\textbf{A.} 人与自然是同一的，自然能够自发满足人的需求 & \textbf{B.} 生产力包括自然要素，合理开发自然能促进社会发展\\[.35em]
\textbf{C.} 人是自然的一部分，人的发展只有适应自然的变化 & \textbf{D.} 人能动改造自然，自然能动补偿人的劳动\\[.35em]\end{tabular}\end{center}

答案：B


3. 马克思在《资本论》中指出: “一个商品占有者出售他现有的商品，而另一个商品占有者却只是作为货币的代表或作为未来货币的代表来购买这种商品。卖者成为债权人，买者成为债务人。由于商品的形态变化或商品的价值形式的发展在这里起了变化货币也就取得了另一种职能。”这里所论述的货币“另一种职能”指的是


\begin{center}\begin{tabular}{>{\raggedright\arraybackslash}p{\dimexpr(\linewidth-8\tabcolsep)/4\relax}>{\raggedright\arraybackslash}p{\dimexpr(\linewidth-8\tabcolsep)/4\relax}>{\raggedright\arraybackslash}p{\dimexpr(\linewidth-8\tabcolsep)/4\relax}>{\raggedright\arraybackslash}p{\dimexpr(\linewidth-8\tabcolsep)/4\relax}}
\textbf{A.} 价值尺度 & \textbf{B.} 流通手段 & \textbf{C.} 贮藏手段 & \textbf{D.} 支付手段\\[.35em]\end{tabular}\end{center}

答案：D


4. 与第二次世界大战前的资本主义相比，当代资本主义在许多方面已经并正在发生着深刻的变化。当代资本主义社会，大公司经营活动的实际控制者是


\begin{center}\begin{tabular}{>{\raggedright\arraybackslash}p{\dimexpr(\linewidth-4\tabcolsep)/2\relax}>{\raggedright\arraybackslash}p{\dimexpr(\linewidth-4\tabcolsep)/2\relax}}
\textbf{A.} 高级职业经理 & \textbf{B.} 股东大会\\[.35em]\end{tabular}\end{center}

\begin{center}\begin{tabular}{>{\raggedright\arraybackslash}p{\dimexpr(\linewidth-4\tabcolsep)/2\relax}>{\raggedright\arraybackslash}p{\dimexpr(\linewidth-4\tabcolsep)/2\relax}}
\textbf{C.} 监事会 & \textbf{D.} 董事长\\[.35em]\end{tabular}\end{center}

答案：A


5. 共同富裕是社会主义的本质要求，是中国式现代化的重要特征。党的十九届五中全会对扎实推动共同富裕作出重大战略部署，明确提出了推进共同富裕的2035年远景目标。这一目标是


\begin{center}\begin{tabular}{>{\raggedright\arraybackslash}p{\dimexpr(\linewidth-4\tabcolsep)/2\relax}>{\raggedright\arraybackslash}p{\dimexpr(\linewidth-4\tabcolsep)/2\relax}}
\textbf{A.} 全体人民共同富裕基本实现 & \textbf{B.} 全体人民共同富裕取得更为明显的实质性进展\\[.35em]
\textbf{C.} 全体人民共同富裕全面实现 & \textbf{D.} 全体人民共同富裕同步实现\\[.35em]\end{tabular}\end{center}

答案：B


6. 发展格局是经济现代化的路径选择，是关系我国发展全局的重大战略任务。立足新发展阶段，贯彻新发展理念，要致力构建以国内大循环为主体、国内国际双循环相互促进的新发展格局。构建新发展格局的关键在于


\begin{center}\begin{tabular}{>{\raggedright\arraybackslash}p{\dimexpr(\linewidth-4\tabcolsep)/2\relax}>{\raggedright\arraybackslash}p{\dimexpr(\linewidth-4\tabcolsep)/2\relax}}
\textbf{A.} 经济循环的畅通无阻 & \textbf{B.} 市场主体的活力\\[.35em]
\textbf{C.} 高水平的自立自强 & \textbf{D.} 产业链供应链的优化升级\\[.35em]\end{tabular}\end{center}

答案：A


7. 解决好“三农”问题始终是全党工作的重中之重。新时代脱贫攻坚目标任务完成以后,“三农” 工作重心将历史性地转向全面推进乡村振兴，2021年的中央一号文件对全面推进乡村振兴，加快农业农村现代化作了进一步的部署。根据这一部署，“十四五”时期“三农” 工作摆在首要位置的是


\begin{center}\begin{tabular}{>{\raggedright\arraybackslash}p{\dimexpr(\linewidth-4\tabcolsep)/2\relax}>{\raggedright\arraybackslash}p{\dimexpr(\linewidth-4\tabcolsep)/2\relax}}
\textbf{A.} 依托乡村特色优势资源，打造农业全产业链 & \textbf{B.} 扩大农村需求，畅通城乡经济循环\\[.35em]
\textbf{C.} 补齐农业农村短板，带动城乡协调发展 & \textbf{D.} 巩固拓展脱贫攻坚成果，守住防止规模性返贫底线\\[.35em]\end{tabular}\end{center}

答案：D


8. 文化软实力集中体现了一个国家基于文化而具有的凝聚力和生命力，以及由此产生的吸引力和影响力。提高国家文化软实力需要强大的国际传播能力，我国加强国际传播能力的建设的重要任务是


\begin{center}\begin{tabular}{>{\raggedright\arraybackslash}p{\dimexpr(\linewidth-4\tabcolsep)/2\relax}>{\raggedright\arraybackslash}p{\dimexpr(\linewidth-4\tabcolsep)/2\relax}}
\textbf{A.} 加快发展新型文化业态 & \textbf{B.} 推进中华传统文化创造性转化，创新性发展\\[.35em]\end{tabular}\end{center}

\begin{center}\begin{tabular}{>{\raggedright\arraybackslash}p{\dimexpr(\linewidth-4\tabcolsep)/2\relax}>{\raggedright\arraybackslash}p{\dimexpr(\linewidth-4\tabcolsep)/2\relax}}
\textbf{C.} 推进社会主义文化深入人心，激发全民族文化创造活力 & \textbf{D.} 讲好中国故事，传播好中国声音，展示真实立体、全面的中国\\[.35em]\end{tabular}\end{center}

答案：D


9. 戊戌维新运动是一次爱国救亡运动，也是一次资产阶级性质的政治改良运动。维新派尽管不敢从根本上否定封建主义，但突破了洋务派“中体西用”思想的局限，主张


\begin{center}\begin{tabular}{>{\raggedright\arraybackslash}p{\dimexpr(\linewidth-4\tabcolsep)/2\relax}>{\raggedright\arraybackslash}p{\dimexpr(\linewidth-4\tabcolsep)/2\relax}}
\textbf{A.} 用君主立宪制取代君主专制制度 & \textbf{B.} 用民主共和制度取代封建专制制度\\[.35em]
\textbf{C.} 用总统共和制取代君主立宪制 & \textbf{D.} 用议会共和制取代君主立宪制\\[.35em]\end{tabular}\end{center}

答案：A


10. 大革命失败后, 要不要坚持革命?如何坚持革命?这是摆在中国共产党面前的两个带根本性的问题。党从残酷的现实中认识到，没有革命的武装就无法战胜武装的反革命，就无法夺取中国革命胜利，就无法改变中国人民和中华民族的命运，必须以武装的革命反对武装的反革命。南昌起义打响了武装反抗国民党反动派的第一枪，这是中国共产党


\begin{center}\begin{tabular}{>{\raggedright\arraybackslash}p{\dimexpr(\linewidth-4\tabcolsep)/2\relax}>{\raggedright\arraybackslash}p{\dimexpr(\linewidth-4\tabcolsep)/2\relax}}
\textbf{A.} 实施土地革命和武装起义方针的开始 & \textbf{B.} 建设无产阶级领导的新型人民军队的开端\\[.35em]
\textbf{C.} 实行工农武装割据的开始 & \textbf{D.} 独立领导革命战争、创建人民军队和武装夺取政权的开端\\[.35em]\end{tabular}\end{center}

答案：D


11. 1949 年 3月 23 日上午，中共中央离开西柏坡向北平进发。临行前，毛泽东把进北平比作“ 进京赶考”，说“我们决不当李自成，我们都希望考个好成绩。”3 月 25 日，毛泽东等中共中央领导人与中央机关、人民解放军总部进驻北平香山，标志着


\begin{center}\begin{tabular}{>{\raggedright\arraybackslash}p{\dimexpr(\linewidth-4\tabcolsep)/2\relax}>{\raggedright\arraybackslash}p{\dimexpr(\linewidth-4\tabcolsep)/2\relax}}
\textbf{A.} 旧民主主义革命向新民主主义革命转变 & \textbf{B.} 新民主主义社会开始向社会主义社会过渡\\[.35em]
\textbf{C.} 中国革命重心从农村转向城市 & \textbf{D.} 民族民主革命任务的完成\\[.35em]\end{tabular}\end{center}

答案：C


12. 1956 年 4月，毛泽东提出把马克思列宁主义基本原理同中国具体实际进行“第二次结合”，其目的是


\begin{center}\begin{tabular}{>{\raggedright\arraybackslash}p{\dimexpr(\linewidth-4\tabcolsep)/2\relax}>{\raggedright\arraybackslash}p{\dimexpr(\linewidth-4\tabcolsep)/2\relax}}
\textbf{A.} 加强中国共产党的自身建设 & \textbf{B.} 为社会主义改造创造条件\\[.35em]
\textbf{C.} 找出在中国怎样建设社会主义的道路 & \textbf{D.} 全面建设社会主义现代化强国\\[.35em]\end{tabular}\end{center}

答案：C


13. 中华传统美德是中华优秀文化的重要组成部分，其内容博大精深、源远流长。从《诗经》中的“夙夜在公”到《尚书》中的“以公灭私”，从西汉贾谊《治安策》中的“国而忘家，公而忘私”到宋代范仲淹《岳阳楼记》中的“先天下之忧而忧，后天下之乐而乐”， 再到清代林则徐的“ 苟利国家生死以，岂因祸福避趋之”，贯穿其中的传统美德是


\begin{center}\begin{tabular}{>{\raggedright\arraybackslash}p{\dimexpr(\linewidth-4\tabcolsep)/2\relax}>{\raggedright\arraybackslash}p{\dimexpr(\linewidth-4\tabcolsep)/2\relax}}
\textbf{A.} 强调知行合一，注重躬行实践 & \textbf{B.} 推崇“仁爱”原则，注重以和为贵\\[.35em]
\textbf{C.} 重视整体利益，强调责任奉献 & \textbf{D.} 提倡人伦价值，重视道德义务\\[.35em]\end{tabular}\end{center}

答案：C


14. 法律作为上层建筑的重要组成部分，不是凭空产生的，也不是永恒存在的，而是由一定的社会物质生活条件所决定的。决定法律本质、内容和发展方向的根本因素是


\begin{center}\begin{tabular}{>{\raggedright\arraybackslash}p{\dimexpr(\linewidth-8\tabcolsep)/4\relax}>{\raggedright\arraybackslash}p{\dimexpr(\linewidth-8\tabcolsep)/4\relax}>{\raggedright\arraybackslash}p{\dimexpr(\linewidth-8\tabcolsep)/4\relax}>{\raggedright\arraybackslash}p{\dimexpr(\linewidth-8\tabcolsep)/4\relax}}
\textbf{A.} 地理环境 & \textbf{B.} 物质资料的生产方式 & \textbf{C.} 人口素质 & \textbf{D.} 统治阶级的意志\\[.35em]\end{tabular}\end{center}

答案：B


15. 2021年9月 27 日至28 日，中央人才工作会议在北京举行。中共中央总书记、国家主席、中央军委主席习近平出席会议并发表重要讲话，强调要坚持党管人才，坚持面向世界科技前沿、面向经济主战场、面向国家重大需求、面向人民生命健康，深入实施新时代人才强国战略，全方位培养、引进、用好人才，加快建设


\begin{center}\begin{tabular}{>{\raggedright\arraybackslash}p{\dimexpr(\linewidth-4\tabcolsep)/2\relax}>{\raggedright\arraybackslash}p{\dimexpr(\linewidth-4\tabcolsep)/2\relax}}
\textbf{A.} 世界重要人才中心和创新高地 & \textbf{B.} 北京、上海、粤港澳大湾区高水平人才高地\\[.35em]
\textbf{C.} 一批国家实验室和新型研发机构 & \textbf{D.} 中心城市吸引和聚集人才的平台\\[.35em]\end{tabular}\end{center}

答案：A


16. 《中华人民共和国反外国制裁法》于 2021 年 6月 10 日起施行。这是一部指向性、针对性颇强的专门法律，共有 16 条。这部法律主要针对的是


\begin{center}\begin{tabular}{>{\raggedright\arraybackslash}p{\dimexpr(\linewidth-4\tabcolsep)/2\relax}>{\raggedright\arraybackslash}p{\dimexpr(\linewidth-4\tabcolsep)/2\relax}}
\textbf{A.} 西方某些大国近年来对我国的“贸易战” & \textbf{B.} 少数国家操纵国际组织挑起的“货币战”\\[.35em]
\textbf{C.} 资本主义大国对社会主义国家的“新冷战” & \textbf{D.} 外国干涉中国内政的所谓“单边制裁”\\[.35em]\end{tabular}\end{center}

答案：D


\subsection*{二、多项选择题：17～33小题，每小题2分，共34分。下列每题给出的四个选项中，至}

少有两个选项是符合题目要求的。请在答题卡上将所选项的字母涂黑。多选或少选均不得分。


17. 在深入推动黄河流域生态保护和高质量发展座谈会上，习近平总书记谈及水资源和发展的关系时，以传统名吃“羊肉泡馍”作形象比喻，强调要全方位贯彻“四水四定”(以水定城、以水定地、以水定人、以水定产)原则, 精打细算用好水资源,“有多少汤泡多少馍” 让水资源用在最该用的地方。“有多少汤泡多少馍”蕴含的哲学道理是


\begin{center}\begin{tabular}{>{\raggedright\arraybackslash}p{\dimexpr(\linewidth-4\tabcolsep)/2\relax}>{\raggedright\arraybackslash}p{\dimexpr(\linewidth-4\tabcolsep)/2\relax}}
\textbf{A.} 创造条件，充分发挥意识能动性 & \textbf{B.} 一切从实际出发，实事求是\\[.35em]
\textbf{C.} 因地制宜，因时制宜 & \textbf{D.} 尊重规律，把握适度原则\\[.35em]\end{tabular}\end{center}

答案：ABCD


18. 人类历史上的每一次科技革命都与材料的发展息息相关，而新材料的研制却是颇为不易的。人工智能可以借助数据共享，对先进材料的物理化学性质进行预测筛选，从而加快新材料的合成和生产。作为人工智能的一个分支，机器学习算法在辅助新材料设计时尤为 “得力”，其工作过程主要包括“描述符”生成、模型构建和验证、材料预测、实验验证等步骤。人工智能辅助新材料研发的过程表明


\begin{center}\begin{tabular}{>{\raggedright\arraybackslash}p{\dimexpr(\linewidth-4\tabcolsep)/2\relax}>{\raggedright\arraybackslash}p{\dimexpr(\linewidth-4\tabcolsep)/2\relax}}
\textbf{A.} 科学研究能够任意改变物质的性质和结构 & \textbf{B.} 人工智能能够取代人类对物质世界的认识\\[.35em]
\textbf{C.} 人类对于物质的认识是一个不断深化的过程 & \textbf{D.} 具体的物质结构和性质的变化并不改变世界的物质性\\[.35em]\end{tabular}\end{center}

答案：CD


19. 时间是万事万物存在的刻度。1秒钟，电影放映24帧画面，猎豹在草原上飞奔28米，蜂鸟振动翅膀55次；1分钟，登山队员攀登珠峰顶峰58.3厘米，“复兴号”前进5833米。时间创造无限可能。有人努力奔跑，用力以赴的冲刺突破极限；有人砥砺前行，以日复一日的坚守辛勤耕耘。原本匀速流动的时间，正是在生生不息的奋斗中，在昂扬奋发的进取中，确定意义、体现价值，进而定义生命的精彩、定格历史的脉动。人们在奋斗中“定义”时间，说明时间是


\begin{center}\begin{tabular}{>{\raggedright\arraybackslash}p{\dimexpr(\linewidth-4\tabcolsep)/2\relax}>{\raggedright\arraybackslash}p{\dimexpr(\linewidth-4\tabcolsep)/2\relax}}
\textbf{A.} 测量事物运动的客观尺度 & \textbf{B.} 物质运动的存在形式\\[.35em]
\textbf{C.} 事物运动的主观联想 & \textbf{D.} 与物质运动不可分割的\\[.35em]\end{tabular}\end{center}

答案：BD


20. 如今，“逆袭”“超燃”“躺平”“凡尔赛”“YYDS”等网络热词令人目不暇接，了解这些网络话的意思用法，几乎成了网上冲浪的必修课。嵌入日常生活交流的网络热词，为语言的发展带来新语料，也能够让人们一窥当下的社会生活和社会心态。网络热词流行的现象表明


\begin{center}\begin{tabular}{>{\raggedright\arraybackslash}p{\dimexpr(\linewidth-4\tabcolsep)/2\relax}>{\raggedright\arraybackslash}p{\dimexpr(\linewidth-4\tabcolsep)/2\relax}}
\textbf{A.} 社会心态是社会生产和生活的反映 & \textbf{B.} 社会意识随着社会存在的发展而变化\\[.35em]
\textbf{C.} 社会意识是主体对社会生活的自发反映 & \textbf{D.} 社会心态的差异性决定社会生活的多样性\\[.35em]\end{tabular}\end{center}

答案：AB


21. 马克思、恩格斯在《共产党宣言》中指出: “资产阶级在它的不到一百年的阶级统治中所创造的生产力，比过去一切世代创造的全部生产力还要多，还要大。自然力的征服， 机器的采用，化学在工业和农业中的应用，轮船的行驶，铁路的通行，电报的使用，整个整个大陆的开垦，河川的通航，仿佛用法术从地下呼唤出来的大量人口——过去哪一个世纪料想到在社会劳动里蕴藏有这样的生产力呢?”资本主义社会生产力迅速发展的主要原因是


\begin{center}\begin{tabular}{>{\raggedright\arraybackslash}p{\dimexpr(\linewidth-4\tabcolsep)/2\relax}>{\raggedright\arraybackslash}p{\dimexpr(\linewidth-4\tabcolsep)/2\relax}}
\textbf{A.} 资本追逐剩余价值的内在动力 & \textbf{B.} 市场激烈竞争的外在压力\\[.35em]
\textbf{C.} 无产阶级反对资产阶级的斗争 & \textbf{D.} 资本主义国家对市场的全面控制\\[.35em]\end{tabular}\end{center}

答案：AB


22. 发展数字经济是把握新一轮科技革命和产业变革新机遇的战略选择。据 2021年8月2日中国信息通信研究院发布的《全球数字经济白皮书》统计显示，2020 年我国数字经济规模近 5.4 万亿美元，居世界第二位:同比增长 9.6\%, 增速位于全球第一。 我国数字经济健康发展有利于


\begin{center}\begin{tabular}{>{\raggedright\arraybackslash}p{\dimexpr(\linewidth-4\tabcolsep)/2\relax}>{\raggedright\arraybackslash}p{\dimexpr(\linewidth-4\tabcolsep)/2\relax}}
\textbf{A.} 推动建设现代化经济体系 & \textbf{B.} 推动构建新发展格局\\[.35em]
\textbf{C.} 推动化解发展中的根本制约因素 & \textbf{D.} 推动构筑国家竞争新优势\\[.35em]\end{tabular}\end{center}

答案：ABD


23. 2021 年是中国加入世界贸易组织 20 周年。20 年来，中国经济总量从世界第六位上升到第二位，货物贸易从世界第六位上升到第一位， 服务贸易从世界第十一位上升到第二位，利用外资稳居发展中国家首位，对外直接投资从世界第二十六位上升到第一位。目前， 中国已成为 50 多个国家和地区的最大贸易伙伴、120 多个国家和地区的前三大贸易伙伴。 中国入世 20 年的历史充分证明


\begin{center}\begin{tabular}{>{\raggedright\arraybackslash}p{\dimexpr(\linewidth-4\tabcolsep)/2\relax}>{\raggedright\arraybackslash}p{\dimexpr(\linewidth-4\tabcolsep)/2\relax}}
\textbf{A.} 我国经济已经实现了由高速增长向高质量发展转变 & \textbf{B.} 以开放促改革、促发展是我国现代化建设的重要法宝\\[.35em]
\textbf{C.} 坚持开放合作才能获得更多发展机遇和更大发展空间 & \textbf{D.} 我国是多边贸易体制的坚定支持者、积极参与者和重要贡献\\[.35em]\end{tabular}\end{center}

答案：BCD


24. 中国共产党自成立以来不断探索和发展适合中国国情的民主道路，使人民民主在东方大国落地生根、繁荣发展。党的十八大以来,我们党不断深化对民主政治发展规律的认识， 践行以人民为中心的发展思想，积极发展全过程人民民主。全过程人民民主


\begin{center}\begin{tabular}{>{\raggedright\arraybackslash}p{\dimexpr(\linewidth-4\tabcolsep)/2\relax}>{\raggedright\arraybackslash}p{\dimexpr(\linewidth-4\tabcolsep)/2\relax}}
\textbf{A.} 是完整制度程序和完整参与实践有机统一的民主 & \textbf{B.} 是全链条、全方位、全覆盖的民主\\[.35em]
\textbf{C.} 是最广泛、最真实、最管用的社会主义民主 & \textbf{D.} 实现了过程民主和成果民主、程序民主和实质民主、直接民主和间接民主、人民民主和国家意志相统一\\[.35em]\end{tabular}\end{center}

答案：ABCD


25. 生态文明的核心是坚持人与自然和谐共生。2021 年 4 月一群栖息在云南西双版纳国家级自然保护区的野生亚洲象，一路向北到达昆明，走过稻田、穿过城镇，历时 110 多天， 行进 1000 多公里。8月上旬，北移的 15 头亚洲象全部安全南返。亚洲象群的迁移引起国内外和社会各界广泛关注。习近平主席在《生物多样性公约》第十五次缔约方大会领导人峰会上的主旨讲话中特别指出: “云南大象的北上及返回之旅，让我们看到了中国保护野生动物的成果。”中国政府一贯高度重视生物多样性保护，不断加强和创新生物多样性保护举措。这些举措主要有


\begin{center}\begin{tabular}{>{\raggedright\arraybackslash}p{\dimexpr(\linewidth-4\tabcolsep)/2\relax}>{\raggedright\arraybackslash}p{\dimexpr(\linewidth-4\tabcolsep)/2\relax}}
\textbf{A.} 划定生态保护红线 & \textbf{B.} 加大自然资源的开发利用\\[.35em]
\textbf{C.} 建立生物安全风险防控和治理体系 & \textbf{D.} 构建以国家公园为主体的自然保护地体系\\[.35em]\end{tabular}\end{center}

答案：ACD


26. 人类各种文明之间的交流是推动世界和平发展的重要动力。习近平主席曾经在比利时布鲁日欧洲学院演讲时指出:“我们要建设文明共荣之桥, 把中欧两大文明连接起来。……正如中国人喜欢茶而比利时人喜爱啤酒一样，茶的含蓄内敛和酒的热烈奔放代表了品味生命、解读世界的两种不同方式。但是，茶和酒并不是不可兼容的，既可以酒逢知己千杯少， 也可以品茶品味品人生。”习近平主席以茶酒巧喻不同文明之间的关系，旨在阐明


\begin{center}\begin{tabular}{>{\raggedright\arraybackslash}p{\dimexpr(\linewidth-4\tabcolsep)/2\relax}>{\raggedright\arraybackslash}p{\dimexpr(\linewidth-4\tabcolsep)/2\relax}}
\textbf{A.} 人类文明多样性是世界的基本特征 & \textbf{B.} 文明因多样而交流，因交流而互鉴，因互鉴而发展\\[.35em]
\textbf{C.} 不同文明凝聚着不同民族的智慧和贡献，没有高低优劣之分 & \textbf{D.} 文明差异不应成为世界冲突的根源，而应成为人类文明进步的动力\\[.35em]\end{tabular}\end{center}

答案：ABCD


27. 1935 年 1 月中共中央政治局在长征途中召开的遵义会议，是党的历史上一个生死攸关


的转折点。这次会议


\begin{center}\begin{tabular}{>{\raggedright\arraybackslash}p{\dimexpr(\linewidth-4\tabcolsep)/2\relax}>{\raggedright\arraybackslash}p{\dimexpr(\linewidth-4\tabcolsep)/2\relax}}
\textbf{A.} 全面讨论了政治路线的是非问题 & \textbf{B.} 解决了当时党内所面临的最迫切的组织问题和军事问题\\[.35em]
\textbf{C.} 结束了“左”倾教条主义错误在中央的统治 & \textbf{D.} 开始确立毛泽东在中共中央和红军的领导地位\\[.35em]\end{tabular}\end{center}

答案：BCD


28. 1981 年 6月党的十一届六中全会通过的《关于建国以来党的若干历史问题的决议》， 科学评价了毛泽东和毛泽东思想的历史地位，把毛泽东思想活的灵魂概括为


\begin{center}\begin{tabular}{>{\raggedright\arraybackslash}p{\dimexpr(\linewidth-8\tabcolsep)/4\relax}>{\raggedright\arraybackslash}p{\dimexpr(\linewidth-8\tabcolsep)/4\relax}>{\raggedright\arraybackslash}p{\dimexpr(\linewidth-8\tabcolsep)/4\relax}>{\raggedright\arraybackslash}p{\dimexpr(\linewidth-8\tabcolsep)/4\relax}}
\textbf{A.} 群众路线 & \textbf{B.} 实事求是 & \textbf{C.} 武装斗争 & \textbf{D.} 独立自主\\[.35em]\end{tabular}\end{center}

答案：ABD


29. 2013 年 11 月，党的十八届三中全会在北京举行。全会审议通过的《中共中央关于全面深化改革若干重大问题的决定》，对全面深化改革作出顶层设计和总体规划。《中共中央关于党的百年奋斗重大成就和历史经验的决议》指出，党的十一届三中全会是划时代的， 党的十八届三中全会也是划时代的。党的十八届三中全会之所以也是划时代的，是因为这次全会


\begin{center}\begin{tabular}{>{\raggedright\arraybackslash}p{\dimexpr(\linewidth-4\tabcolsep)/2\relax}>{\raggedright\arraybackslash}p{\dimexpr(\linewidth-4\tabcolsep)/2\relax}}
\textbf{A.} 实现改革由局部探索、破冰突围到系统集成、全面深化的转变 & \textbf{B.} 开启了改革开放和社会主义现代化建设新时期\\[.35em]
\textbf{C.} 开创了我国改革开放新局面 & \textbf{D.} 确定建立社会主义市场经济体制\\[.35em]\end{tabular}\end{center}

答案：AC


30. 2020 年我国颁布的民法典具有鲜明道德导向。如第 7条“民事主体从事民事活动，应当遵循诚信原则，秉持诚实，恪守承诺。”第 8 条“民事主体从事民事活动，不得违反法律，不得违背公序良俗。”第 184 条“因自愿实施紧急救助行为造成受助人损害的，救助人不承担民事责任。”法律对道德的导向主要表现为


\begin{center}\begin{tabular}{>{\raggedright\arraybackslash}p{\dimexpr(\linewidth-4\tabcolsep)/2\relax}>{\raggedright\arraybackslash}p{\dimexpr(\linewidth-4\tabcolsep)/2\relax}}
\textbf{A.} 法律助推个人品德养成 & \textbf{B.} 法律决定道德传承的效果\\[.35em]
\textbf{C.} 法律助推社会公德水平提升 & \textbf{D.} 法律为弘扬美德提供保障\\[.35em]\end{tabular}\end{center}

答案：ACD


31. 习近平法治思想是全面依法治国的根本遵循和行动指南。2020 年 11 月，习近平在中央全面依法治国工作会议上，用“十一个坚持”对全面依法治国进行了系统阐释、部署，从全面依法治国的政治方向、战略地位、工作布局、主要任务、重大关系、重要保障等方面提出了一系列新理念新观点新论断。其中关于全面依法治国政治方向的是


\begin{center}\begin{tabular}{>{\raggedright\arraybackslash}p{\dimexpr(\linewidth-4\tabcolsep)/2\relax}>{\raggedright\arraybackslash}p{\dimexpr(\linewidth-4\tabcolsep)/2\relax}}
\textbf{A.} 坚持党对全面依法治国的领导 & \textbf{B.} 坚持以人民为中心\\[.35em]
\textbf{C.} 坚持中国特色社会主义法治道路 & \textbf{D.} 坚持统筹推进国内法治和涉外法治\\[.35em]\end{tabular}\end{center}

答案：ABC


32. 2021年5月15日7时18分，天问一号探测器成功着陆于火星乌托邦平原南部预选着陆区，我国首次火星探测任务着陆火星取得成功。其重要意义表现在


\begin{center}\begin{tabular}{>{\raggedright\arraybackslash}p{\dimexpr(\linewidth-4\tabcolsep)/2\relax}>{\raggedright\arraybackslash}p{\dimexpr(\linewidth-4\tabcolsep)/2\relax}}
\textbf{A.} 我国在行星探测领域进入世界先进行列 & \textbf{B.} 我国首次自主执行了火星探测任务\\[.35em]
\textbf{C.} 为探索宇宙奥秘、增进对火星演化的认知、了解生命起源作出了中国贡献 & \textbf{D.} 我国航天事业发展取得了又一具有里程碑意义的进展\\[.35em]\end{tabular}\end{center}

答案：ABCD


33. 2021年11月16日，中国国家主席习近平同美国总统拜登举行视频会晤，双方就事关中美关系发展的战略性、全局性、根本性问题以及共同关心的其他重要问题进行了充分、 深入的沟通和交流。习近平强调，新时期中美相处应该坚持的原则是


\begin{center}\begin{tabular}{>{\raggedright\arraybackslash}p{\dimexpr(\linewidth-8\tabcolsep)/4\relax}>{\raggedright\arraybackslash}p{\dimexpr(\linewidth-8\tabcolsep)/4\relax}>{\raggedright\arraybackslash}p{\dimexpr(\linewidth-8\tabcolsep)/4\relax}>{\raggedright\arraybackslash}p{\dimexpr(\linewidth-8\tabcolsep)/4\relax}}
\textbf{A.} 相互尊重 & \textbf{B.} 和平共处 & \textbf{C.} 合作共赢 & \textbf{D.} 结伴不结盟\\[.35em]\end{tabular}\end{center}

答案：ABC


\subsection*{三、分析题：34～38小题，每小题10分，共50分。要求结合所学知识分析材料回答问题。将答案写在答题卡指定位置的边框区域内。}

34. 结合材料回答问题:

材料 1

进入新时代，我国面临更为严峻的国家安全形势，外部压力前所未有，传统安全威胁和非传统安全威胁相互交织，“黑天鹅”、“灰犀牛”事件时有发生。同形势任务要求相比，我国维护国家安全能力不足，应对各种重大风险能力不强，维护国家安全的统筹协调机制不健全。党中央强调，国泰民安是人民群众最基本、最普遍的愿望。必须坚持底线思

维、居安思危、未雨绸缪，坚持国家利益至上，以人民安全为宗旨，以政治安全为根本， 以经济安全为基础，以军事、科技、文化、社会安全为保障，以促进国际安全为依托，统筹发展和安全,统筹开放和安全，统筹传统安全和非传统安全, 统筹自身安全和共同安全， 统筹维护国家安全和塑造国家安全。

摘自《中共中央关于党的百年奋斗重大成就和历史经验的决议》

材料 2

湖北“黑天鹅”事件通常指现实生活中出现的出乎预料的小概率事件。在特定时间内发生的可能性相对较低的小概率事件，广泛存在于自然、经济、政治等各个领域，具有偶发性、难以预测性等特征。小概率事件虽然发生概率小，但并非零概率事件，若从长时段来看，只要具备相关因素和条件，就可能会发生。小概率事件的影响不局限于一时一地， 一旦发生，就可能会形成多米诺骨牌效应，导致系统性风险，给整个人类社会发展带来深远影响。

习近平总书记在庆祝中国共产党成立 100 周年大会上的讲话中指出，新的征程上，我们必须增强忧患意识、始终居安思危,深刻认识我国社会主要矛盾变化带来的新特征新要求。 深刻认识错综复杂的国际环境带来的新机遇新挑战，敢于斗争，善于斗争，逢山开道、过水架桥，勇于战胜一切风险挑战。应对我国发展环境深刻复杂变化，特别是其中隐藏的重大风险挑战，是向着全面建成社会主义现代化强国目标迈进的必然要求，我国发展具有多方面优势和条件，但我国发展不平衡不充分问题仍然突出，重难点领域关键环节改革任务仍然艰巨。发展环境的深刻复杂变化，既要求我们牢牢把握机遇发展自己，又要求我们树立辩证思维。提高驾驭复杂局面、处理复杂问题的本领。

摘编自《人民日报》（2020年7月27日、2021年7月2日）

（1） 运用必然与偶然辩证关系的原理，说明小概率事件并非零概率事件。 （5分）

（2） 面对复杂局面应如何运用好辩证思维? （5分）

答题思路:

(1) 第一问考查必然和偶然的关系，首先论述必然和偶然的关系，其次结合材料说明小概率事件也可能导致系统性风险。

(2) 第二问考查辩证思维，辩证思维就是坚持联系、发展和矛盾的观点。可以分别论述联系、发展和矛盾的相关原理，其次结合材料说明坚持辩证思维的必要性，尤其是矛盾分析法在材料中体现的比较突出。

35. 结合材料回答问题;

在中国共产党成立 100 周年的重要历史时刻,在党和人民胜利实现第一个百年奋斗目标全面建成小康社会，正在向着全面建成社会主义现代化强国的第二个百年奋斗目标迈进

的重大历史关头，党的十九届六中全会于2021年11月8日至11日在北京胜利举行，审议通过了《中共中央关于党的百年奋斗重大成就和历史经验的决议》。全会聚焦总结党的百年奋斗重大成就和历史经验，深入研究我们党不断推进马克思主义中国化的百年历程， 把“坚持理论创新”概括为党百年奋斗的十条历史经验之一，深刻指出党的百年奋斗展示了马克思主义的强大生命力。

“我们党的历史，就是一部不断推进马克思主义中国化的历史，就是一部不断推进理论创新、进行理论创造的历史。”一百年来，中国共产党人“顶”马克思主义的“天”， “立”中国国情的“地”，在马克思主义及其中国化创新理论的指引下，团结带领中国人民，书写了中华民族几千年历史上最恢宏的史诗。实践充分证明，党之所以能够领导人民在一次次求索、一次次挫折、一次次开拓中完成中国其他各种政治力量不可能完成的艰巨任务，根本在于坚持解放思想、实事求是、与时俱进、求真务实，坚持把马克思主义基本原理同中国具体实际相结合、同中华优秀传统文化相结合，坚持实践是检验真理的唯一标准，坚持一切从实际出发。及时回答时代之问、人民之问，不断推进马克思主义中国化时代化，用马克思主义中国化的科学理论引领伟大实践。

中国特色社会主义进入新时代。中国共产党面临的主要任务是，实现第一个百年奋斗目标，开启实现第二个百年奋斗目标新征程，朝着实现中华民族伟大复兴的宏伟目标继续前进。在这个风云激荡、波澜壮阔的新时代，世情国情党情持续发生深刻复杂变化，向中国共产党人提出一系列新的重大课题。《决议》突出中国特色社会主义新时代这个重点，在党的十九大报告“八个明确”的基础上，用“十个明确”对习近平新时代中国特色社会主义思想的核心内容作了进一步概括和阐述，《决议》还从 13 个方面分领域总结新时代党领导人民取得的巨大成就，并重点概括了其中的原创性的理念和思想。这些战略思想和创新理念，是党对中国特色社会主义建设规律认识深化和理论创新的重大成果。《决议》指出，习近平新时代中国特色社会主义思想实现了马克思主义中国化新的飞跃。党确立习近平同志党中央的核心，全党的核心地位，确立习近平新时代中国特色社会主义思想的指导地位，反映了全党全军全国各族人民共同心愿，对新时代党和国家事业发展、对推进中华民族伟大复兴历史进程具有决定性意义。

摘编自《求是》（2021年第22期）、《人民日报》（2021年11月16日）

（1） 怎样认识我们党的历史“就是一部不断推进理论创新、进行理论创造的历史”? （5分）

(2) 结合中华民族伟大复兴战略全局和世界百年未有之大变局，阐述如何理解习近平新时代中国特色社会主义思想“实现了马克思主义中国化新的飞跃”。 （5分）

答题思路:

(1) 第一问回答马克思主义中国化的三次飞跃成果: 第一大成果是毛泽东思想，第二大成果是中特理论体系; 第三大成果是习近平新时代中国特色社会主义思想。

(2) 第一，这一思想深刻回答了新时代的重大时代课题; 第二，这一思想提出了一系列原创性的治国理政新理念新思想新战略; 第三，这一思想是中华文化和中国精神的时代精华。

总体来讲，习近平新时代中国特色社会主义思想，是把马克思主义基本原理同中国具体实际相结合、同中华优秀传统文化相结合的重大理论成果，是党的十八大以来历史性成就和历史性变革的重要理论结晶，实现了马克思主义中国化的历史性飞跃和创造性升华。

36. 结合材料回答问题:

材料 1

2016年11月11 日,习近平总书记在纪念孙中山先生诞辰 150 周年大会上的讲话中指出:

“孙中山先生在从事紧张的革命活动的过程中, 一直思考着建设中国的问题。1917 年到 1919 年，他写出《建国方略》一书，构想了中国建设的宏伟蓝图，其中提出要修建约 16 万公里的铁路，把中国沿海、内地、边疆连接起来; 修建 160 万公里的公路，形成遍布全国的公路网，并进入青藏高原，开凿和整修全国水道和运河，建设三峡大坝，发展内河交通和水利、电力事业; 在中国北部、中部、南部、沿海各修建一个世界水平的大海港; 大力发展农业、制造业、矿业等等，孙中山先生擘画的这个蓝图，显示了他对中国发展的卓越见解和强烈期盼。当时，有的外国记者认为孙中山先生的这些设想完全是一种空想， 是不可能实现的。”

“的确, 在旧中国的政治经济社会条件下, 孙中山先生的这些宏大构想是难以实现的。 今天，在中国共产党领导下，在全国各族人民顽强奋斗下，孙中山先生当年描绘的这个蓝图早已实现，中国人民创造的许多成就远远超出了孙中山先生的设想。祖国大地上，铁路进青藏，公路密成网，高峡出平湖，港口连五洋， 产业门类齐，稻麦遍地香，神舟遨太空， 国防更坚强。 孙中山先生致力于建设的独立, 民主、 富强的国家早已巍然屹立在世界东方。”

摘自习近平《在纪念孙中山先生诞辰 150 周年大会上的讲话》

材料 2

2021年7月 1 日, 庆祝中国共产党成立 100 周年大会在北京天安门广场隆重举行，习近平总书记在庆祝大会上发表重要讲话，深情回顾中国共产党百年奋斗的光辉历程，高度评价一百年来中国共产党团结带领中国人民创造的伟大成就。习近平总书记强调: “一百年来，中国共产党团结带领中国人民进行的一切奋斗、一切牺牲、一切创造，归结起来就

是一个主题: 实现中华民族伟大复兴。”习近平总书记指出: “一百年，中国共产党的先驱们创建了中国共产党，形成了坚持真理，坚守理想，践行初心、担当使命，不怕牺牲， 英勇斗争，对党忠诚、不负人民的伟大建党精神。这是中国共产党的精神之源。”

回望百年奋斗历程，我们浴血奋战、百折不挠，创造了新民主主义革命的伟大成就， 以英勇顽强的奋斗向世界庄严宣告，中国人民站起来了，中华民族任人宰割、饱受欺凌的时代一去不复返了! 我们自力更生、发愤图强，创造了社会主义革命和建设的伟大成就， 以英勇顽强的奋斗向世界庄严宣告，中国人民不但善于破坏一个旧世界、也善于建设一个新世界，只有社会主义才能救中国，只有社会主义才能发展中国。 我们解放思想、锐意进取，创造了改革开放和社会主义现代化建设的伟大成就，以英勇顽强的奋斗向世界庄严宣告，改革开放是决定当代中国前途命运的关键一招，中国大踏步赶上了时代。我们自信自强，守正创新，统揽伟大斗争、伟大工程、伟大事业、伟大梦想，创造了新时代中国特色社会主义的伟大成就，以英勇顽强的奋斗向世界庄严宣告，中华民族迎来了从站起来、富起来到强起来的伟大飞跃，实现了中华民族伟大复兴，进入了不可逆转的历史进程! 这一百年来开辟的伟大道路，创造的伟大事业、取得的伟大成就， 必将载入中华民族发展史册，人类文明发展史册!

摘编自《人民日报》（2021年7月2日）

（1） 为什么“在旧中国的政治经济社会条件下”，孙中山振兴中华的宏大构想“难以实现”? (5分)

(2) 中国共产党百年奋斗的历史意义是什么? （5分）

答题思路:

(1) 主要是考察辛亥革命失败的原因。最主要的原因就是因为没有坚强的领导核心，没有科学的理论指导，最终没有完成反帝反封建的历史使命.

(2) 第一，从根本上改变了中国人民的前途命运。 第二，开辟了实现中华民族伟大复兴的正确道路。 第三，展示了马克思主义的强大生命力。 第四，深刻影响了世界历史进程。 第五，锻造了走在时代前列的中国共产党。

37. 结合材料回答问题:

材料 1

2021 年 6月 29 日，中共中央将党内最高荣誉“七一勋章”追授黄文秀同志。黄文秀曾被追授“全国脱贫攻坚楷模”“全国优秀共产党员”“时代楷模”等荣誉称号。

黄文秀生在百色、长在百色，学成归来回到百色建设家乡。2018 年 3 月，她主动要求到条

件更为艰苦的百坭村担任驻村第一书记。驻村期间，黄文秀带领村干部和群众,埋头苦干， 学经验、找路子、筹资金，硬化通屯路，安装路灯，修建蓄水池，发展经济林木和砂糖橘种植业，百坭村贫困发生率从 22.8\%降至2.71\%。2019 年，百坭村实现整村脱贫。

在服务百坭村的日子里，黄文秀访遍全部贫困户，帮助村民发展扶贫产业，以真抓实干的作风赢得村民信任;她全身心扑在扶贫事业上，在城乡、村屯间穿梭，用真情奉献与群众打成一片，以使命担当兑现着“不获全胜，决不收兵”的驻村诺言。

作为一名党员干部，黄文秀始终扎根群众、心系群众。作为驻村第一书记，黄文秀还特别注重在脱贫攻坚中发挥基层党支部的战斗堡垒作用。黄文秀始终牢记当年在入党申请书中写下的誓言: “一个人要活得有意义，生存得有价值，就不能光为自己而活，要用自己的力量为他人、为国家、为民族、为社会做出贡献。”

2019 年 6 月，百坭村连降暴雨。因惦记百坭村的防汛抗洪工作,利用周末看望手术不久的父亲后，黄文秀冒雨连夜返回，不料途中遭遇山洪，年轻的生命定格在 30 岁。习近平总书记在对黄文秀先进事迹的指示中强调，黄文秀同志“在脱贫攻坚第一线倾情投入、 奉献自我，用美好青春诠释了共产党人的初心使命，谱写了新时代的青春之歌。”

“年轻人的态度就是乡村的未来”，黄文秀生前曾在驻村日记中这样写道。如今，她的先进事迹和优秀品质被广为传颂，激励着越来越多的青年扎根新农村，创造更美好的生活。

摘编自《人民日报》（2021年6月18日、7月2日、7月27日）

材料 2

一代人有一代人的长征, 一代人有一代人的担当。 每一代青年都有自己的际遇和机缘， 都要在自己所处的时代条件下谋划人生、创造历史。“清澈的爱，只为中国”，这是年轻边防战士在皑皑雪山立下的铿锵誓言；“我们的征途是星辰大海”，这是年轻航天人探索宇宙抒发的豪情壮志；“请党放心，强国有我”，这是年轻学子发自心底的庄严承诺。

青年向上，国家向前。在庆祝中国共产党成立 100 周年大会上，习近平总书记深情寄语新时代的中国青年，“要增强做中国人的志气、骨气、底气，不负时代，不负韶华，不负党和人民的殷切期望！”

摘编自《人民日报》（2021年8月9日）

（1） 黄文秀的先进事迹启示新时代青年应该有怎样的人生态度? （5分）

(2) 一代人有一代人的担当，新时代中国青年应当承担什么样的历史重任? （5分）

答题思路:

（1） 第一问考查人生观中的人生态度问题，首先指出人生态度是人生观的重要内容，正确的人生态度会对世界观和人生观产生重要影响，其次结合黄文秀的先进事迹指出新时代

青年应该以认真、务实、乐观和进取的态度面对人生。

（2） 第二问考查新时代中国青年的历史使命，需要先回答青年担负着实现伟大复兴中国梦的历史重任, 其次可以从人生价值、理想信念等角度说明青年应如何担负这一历史重任。

38. 结合材料回答问题:

材料 1

1945 年，世界反法西斯战争取得胜利。“欲免后世再遭今代人类两度身历惨不堪言之战祸”，刚刚走出浩劫的人类苦苦求索和平进步之道。同年 10 月 24 日，《联合国宪章》 生效，联合国正式成立。

自1949年10月1日成立那天起，中华人民共和国理所当然地享有中国在联合国的席位。1949年11 月 15日，周恩来致电联合国秘书长赖伊和第四届联合国大会主席罗慕洛， 要求“立即取消‘中国国民政府代表团’继续代表中国人民参加联合国的一切权利”。然而，由于美国等西方国家的阻挠，新中国被长期排除在联合国之外，为争取合法权利，中国政府和人民开展了不懈的、不疲倦的斗争，得到世界上主持正义国家的共同支持，最终汇聚成无法阻挡的时代浪潮。

1971 年, 美国政府虽已认识到很难再把中国拒于联合国大门之外, 但又不甘心失败，便与日本政府再次提出“重要问题”案并炮制所谓“双重代表权”案。对此，中国政府严正声明决不允许在联合国出现“两个中国”或“一中一台”的局面。

1971年10月25日，第二十六届联合国大会进行着一场没有硝烟的交锋。由阿尔巴尼亚、阿尔及利亚等 23 个国家共同提交的“要求恢复中华人民共和国在联合国的一切合法权利、立即把台湾国民党当局的代表从联合国及其所属一切机构中驱逐出去”的决议草案， 遭到美国代表的极力阻挠。广大发展中国家代表据理力争，逐一排除了美国设置的障碍。

最终，大会以76 票赞成、35 票反对、17 票弃权的压倒多数通过第 2758 号决议，决定恢复中华人民共和国在联合国的一切权利，承认中华人民共和国政府的代表是中国在联合国的唯一合法代表。 这是中国人民的胜利, 也是世界各国人民的胜利。美国代表也不得不承认: “任何人都不能回避这样一个事实——刚刚投票的结果，确实代表着大多数联合国会员国的看法。”

摘编自《人民日报》（2021年2月23日、10月26日）、《环球人物》（2021年第21期）

材料 2

今年是中华人民共和国恢复联合国合法席位 50 周年。这是中国践行多边主义的 50 年， 是中国全面参与和支持联合国事业的 50 年，是推动全球治理体系向着更加公正合理的方向发展的 50 年。从“三个世界”划分到和平与发展是时代主题的科学判断，从建设和谐

世界到推动构建人类命运共同体……中国在每个历史关头，都为推动世界和平与发展贡献智慧和方案，并以自身发展书写了人类进步的传奇。

当前，百年变局与世纪疫情相互交织，全球治理体系经历深刻调整。某些国家打着多边主义幌子，大搞零和博弈、霸权主义、集团政治、意识形态小圈子等形形色色的伪多边主义，制造新的分裂、引发新的冲突。

世界需要真正的多边主义。和平、发展、公平、正义、民主、自由是全人类的共同价值。习近平在第七十六届联合国大会一般性辩论上发表重要讲话指出: “世界只有一个体系，就是以联合国为核心的国际体系。只有一个秩序，就是以国际法为基础的国际秩序。 只有一套规则，就是以联合国宪章宗旨和原则为基础的国际关系基本准则。”站在新的历史起点，中国将继续全力参与联合国事务，全力捍卫联合国的地位，全力维护联合国宪章宗旨和原则，同世界上一切进步力量携手，共同践行真正的多边主义，向着构建人类命运共同体不断迈进。

摘编自《人民日报》 (2021年9月22日、10月25日、11月30日)

（1） 为什么说中华人民共和国恢复在联合国的一切权利是中国人民的胜利，也是世界各国人民的胜利? （5分）

(2) 习近平主席所强调的世界只有“一个体系”“一个秩序”“一套规则”，对于推动构建人类命运共同体有何重要意义? （5分）

答题思路:

(1) 第一问主要是从中国是联合国第一个签字的国家，而且是第一个发展中国家的角度去阐发在联合国中的地位，主要是答三个 50 年，既要回答对中国的影响，也要回答对世界的影响。

(2) 第二问主要扣着命运共同体，所以需要从政治上、经济上、安全上、生态上以及文化上五个方面去回答如何构建命运共同体。
\end{document}
